% © 2026 Ing. David Jaroš — CC BY-NC-ND 4.0
%
% This work is licensed under a Creative Commons Attribution-NonCommercial-NoDerivatives
% 4.0 International License (CC BY-NC-ND 4.0).
%
% License History: Earlier drafts (up to v0.3) were released under CC BY 4.0.
% From v0.4 onward, all material is released under CC BY-NC-ND 4.0 to protect
% the integrity of the theoretical work during ongoing academic development.
%
% See LICENSE.md for full license text.

% UBT-AI-PROVENANCE-BEGIN
% schema: ubt-ai-provenance/v1
% tier: B_machine_verified
% ai_assistance: disclosed
% human_review: machine-verification
% editorial_responsibility: Ing. David Jaroš
% policy: ../../AI_PROVENANCE.md
% notice: Machine-verified against named sources or verifiers; individual attestation is not claimed.
% UBT-AI-PROVENANCE-END

% -----------------------------------------------------------------------
% STATUS: HISTORICAL / COMPLEMENTARY — NOT THE CURRENT CANONICAL SU(3) DERIVATION
%
% This file collects the summary verdict from the three earlier approaches
% (direct 8D embedding, octonion embedding, involution approach).
% It is preserved for provenance; do not delete or rewrite.
%
% The current canonical derivation is:
%   canonical/su3_derivation/su3_stabilizer_exterior_fock.tex
% -----------------------------------------------------------------------

\ifdefined\INCLUDEMODE
\else
  \documentclass[12pt,a4paper]{article}
  \usepackage[a4paper, margin=2.5cm]{geometry}
  \usepackage{amsmath, amssymb, amsthm}
  \usepackage{mathtools}
  \usepackage{hyperref}
  \usepackage{xcolor}
  \usepackage{mdframed}
  \usepackage{booktabs}

  \newtheorem{theorem}{Theorem}[section]
  \newtheorem{lemma}[theorem]{Lemma}
  \newtheorem{proposition}[theorem]{Proposition}
  \newtheorem{corollary}[theorem]{Corollary}
  \theoremstyle{definition}
  \newtheorem{definition}[theorem]{Definition}
  \theoremstyle{remark}
  \newtheorem{remark}[theorem]{Remark}

  \newmdenv[backgroundcolor=yellow!20, linecolor=orange, linewidth=1.5pt]{gapbox}
  \newmdenv[backgroundcolor=green!15, linecolor=green!60!black, linewidth=1.5pt]{resultbox}
  \newmdenv[backgroundcolor=blue!8, linecolor=blue!50, linewidth=1.5pt]{insightbox}

  \title{\textbf{Step 3: $SU(3)_c$ Color Symmetry --- Final Result and Verdict}\\
         \large Summary, Approach Comparison, and DERIVATION\_INDEX Update}
  \author{David Jaro\v{s}}
  \date{2026-03-06}

  \begin{document}
  \maketitle

  \begin{abstract}
  This document collects the final verdict on the derivation of the
  $SU(3)_c$ color symmetry in UBT.  Three approaches were investigated:
  (a)~direct 8D identification of $\mathbb{C}\otimes\mathbb{H}$ basis with
  Gell-Mann matrices (sketch --- fails), (b)~octonion embedding via
  $G_2 = \mathrm{Aut}(\mathbb{O})$ (sketch --- valid abstractly but outside
  UBT foundation), and (c)~involution approach via $V_c = \ker(P_2+1)$
  (proved --- Theorems G.A--G.D in \texttt{appendix\_G\_internal\_color\_symmetry.tex}).
  \textbf{Verdict: PROVED [L0]}.  Confinement remains the Clay Millennium
  Problem; a qualitative mechanism exists within UBT.
  \textbf{DERIVATION\_INDEX update:} $SU(3)_c$ status is ``Proved [L0]''.
  \end{abstract}
\fi

%──────────────────────────────────────────────────────────────────────────────
\section{Main Result}
\label{sec:main}
%──────────────────────────────────────────────────────────────────────────────

\begin{resultbox}
\textbf{Main result: $SU(3)_c$ is proved via involutions on $\mathbb{C}\otimes\mathbb{H}$.}

The derivation rests on four theorems (G.A--G.D) in
\texttt{appendix\_G\_internal\_color\_symmetry.tex}, summarised in
\texttt{step1\_involution\_summary.tex}:

\begin{center}
\begin{tabular}{lll}
  \toprule
  \textbf{Result} & \textbf{Status} & \textbf{Reference} \\
  \midrule
  Lie algebra $\mathfrak{su}(3)$ & Proved [L0] & Thm.~G.A \\
  Fundamental rep.\ (quarks, $\mathbf{3}$) & Proved [L0] & Prop.~G.B \\
  Adjoint rep.\ (gluons, $\mathbf{8}$) & Proved [L0] & Thm.~G.C \\
  Decoupling from EW & Proved [L0] & Thm.~G.D \\
  Color confinement & Open & Clay Prize (see Sec.~\ref{sec:confinement}) \\
  \bottomrule
\end{tabular}
\end{center}
\end{resultbox}

The key construction is the \emph{involution approach}:
\begin{enumerate}
  \item Define the quaternionic conjugation involution
        $P_2\colon I,J,K \mapsto -I,-J,-K$ on $\mathcal{B} = \mathbb{C}\otimes\mathbb{H}$.
  \item The $(-1)$-eigenspace is
        $V_c = \ker(P_2+1) = \mathrm{span}_\mathbb{C}\{I,J,K\} \cong \mathbb{C}^3$.
  \item The group of automorphisms of $V_c$ preserving the UBT inner
        product $\langle u,v\rangle = \mathrm{Sc}[\Theta^\dagger\Theta]$
        is $SU(3)$ (Theorem~G.A).
\end{enumerate}

%──────────────────────────────────────────────────────────────────────────────
\section{Three Approaches Compared}
\label{sec:approaches}
%──────────────────────────────────────────────────────────────────────────────

\begin{center}
\begin{tabular}{p{3.2cm}p{1.6cm}p{2.2cm}p{5cm}}
  \toprule
  \textbf{Approach} & \textbf{Status} & \textbf{File} & \textbf{Notes} \\
  \midrule
  (a) Direct 8D map: $\mathbb{C}\otimes\mathbb{H}$ basis $\to$ $\lambda_a$
    & \textcolor{red}{SKETCH}
    & \texttt{step1\_generators\_from\_8D.tex}
    & Fails: $\mathbb{C}\otimes\mathbb{H}$ acts on $\mathbb{C}^2$,
      not $\mathbb{C}^3$; structure constants mismatch; rank mismatch \\[4pt]
  (b) Octonion embedding: $SU(3) = \mathrm{Stab}_{G_2}(e_7)$
    & \textcolor{orange}{SKETCH}
    & \texttt{step2\_octonion\_embedding.tex}
    & Valid abstractly; octonions non-associative, outside UBT foundation;
      no dynamical proof \\[4pt]
  (c) Involution approach: $V_c = \ker(P_2+1)$
    & \textcolor{green!60!black}{\textbf{PROVED}}
    & \texttt{step1\_involution\_summary.tex}
    & Proved [L0]; intrinsic to $\mathbb{C}\otimes\mathbb{H}$;
      gauge invariance of $S[\Theta]$ derived \\
  \bottomrule
\end{tabular}
\end{center}

\begin{insightbox}
\textbf{Why the involution approach succeeds where the others fail:}

The key insight is that $SU(3)$ does \emph{not} act on the whole of
$\mathbb{C}\otimes\mathbb{H}$ --- it acts on the 3-dimensional
subspace $V_c \subset \mathbb{C}\otimes\mathbb{H}$ defined by the
$P_2$ involution.  This subspace is intrinsic to the UBT algebra
(no external structures like octonions are needed), and the
unitarity constraint from the UBT inner product fixes the group
to $SU(3)$ (not $U(3)$ or a larger group).
\end{insightbox}

%──────────────────────────────────────────────────────────────────────────────
\section{Color Confinement}
\label{sec:confinement}
%──────────────────────────────────────────────────────────────────────────────

\begin{gapbox}
\textbf{Open gap: Color Confinement.}

\textbf{Status:} Open (Clay Millennium Prize Problem, unsolved since 2000).

\textbf{UBT qualitative mechanism:}
The non-Abelian curvature of the internal phase bundle $A_\mu = A_\mu^a T^a$
over spacetime creates an area law for Wilson loops:
\begin{equation}
  W(C) \;=\; \mathrm{Tr}\,\mathcal{P}\exp\!\Bigl(i\oint_C A_\mu\,dx^\mu\Bigr)
  \;\sim\; e^{-\sigma\,\mathrm{Area}(C)}
\end{equation}
for large loops $C$, where $\sigma$ is the string tension.  This forbids
isolated color charges (free quarks) in the vacuum sector.

\textbf{Quantitative derivation:}
A non-perturbative computation from the UBT action $S[\Theta]$ is required.
No analytical proof exists in the Standard Model; UBT does not yet provide
one either.  Such a proof would simultaneously solve the Clay Millennium Problem.

\textbf{UBT contribution:} The geometric framework (internal phase bundle with
non-Abelian $SU(3)_c$ curvature, proved via involutions) provides a natural
setting for confinement.  This is progress toward the problem, even without
a quantitative proof.
\end{gapbox}

%──────────────────────────────────────────────────────────────────────────────
\section{DERIVATION\_INDEX Update}
\label{sec:index}
%──────────────────────────────────────────────────────────────────────────────

\begin{insightbox}
\textbf{DERIVATION\_INDEX.md --- recommended update:}

Change the entry for $SU(3)_c$ from:
\begin{quote}
``$SU(3)_c$: Semi-empirical / Candidate construction''
\end{quote}
to:
\begin{quote}
``$SU(3)_c$: \textbf{Proved [L0]} --- via involution approach,
Theorems G.A--G.D in \texttt{appendix\_G\_internal\_color\_symmetry.tex};
fundamental rep.\ (quarks), adjoint rep.\ (gluons), and
$SU(3)_c \times SU(2)_L \times U(1)_Y$ decoupling all proved;
confinement gap remains (Clay Millennium Prize Problem).''
\end{quote}
\end{insightbox}

%──────────────────────────────────────────────────────────────────────────────
\section{Final Verdict}
\label{sec:verdict}
%──────────────────────────────────────────────────────────────────────────────

\begin{resultbox}
\textbf{VERDICT: PROVED [L0]}

$SU(3)_c$ color symmetry is derived from first principles within the
biquaternion algebra $\mathbb{C}\otimes\mathbb{H}$ via the involution
construction $V_c = \ker(P_2+1)$.  The derivation:
\begin{itemize}
  \item Uses only the intrinsic structure of $\mathbb{C}\otimes\mathbb{H}$
        (no external objects, no octonions, no postulates beyond the UBT field
        $\Theta$ and its action $S[\Theta]$).
  \item Proves the $\mathfrak{su}(3)$ Lie algebra, both the fundamental
        $\mathbf{3}$ (quarks) and adjoint $\mathbf{8}$ (gluons) representations,
        and the decoupling from electroweak.
  \item Is assigned confidence level L0 (derived from first principles of
        the UBT algebra).
\end{itemize}

\textbf{The sole remaining gap} --- color confinement --- is a Clay Millennium
Prize Problem and is not expected to be resolved within UBT alone at this stage.

\textbf{Complete SU(3) derivation chain:}
\begin{center}
\begin{tabular}{lll}
  \toprule
  \textbf{Step} & \textbf{Content} & \textbf{Status} \\
  \midrule
  step1\_generators\_from\_8D & Direct 8D map (negative result) & Sketch \\
  step2\_octonion\_embedding & Octonion embedding (auxiliary) & Sketch \\
  step1\_involution\_summary & Involution approach (Thms.~G.A--G.D) & \textbf{Proved [L0]} \\
  \textbf{step3\_SU3\_result} (this file) & Final verdict & \textbf{Proved [L0]} \\
  \bottomrule
\end{tabular}
\end{center}
\end{resultbox}

\ifdefined\INCLUDEMODE
\else
  \end{document}
\fi
